\subsection{Bounded development study: interventions and selection}

The protocol was frozen before the 34-run study on 4 November 2026. It reused the 144/49 partition and the historical 168-slide external collection. Ten configurations were trained at seeds 17, 29 and 43. The two controls repeated the full learned-gate Hybrid and matched-attention configurations above. Their reference scores, selected epochs and decisions were checked against the earlier controlled runs, with score tolerance $3\times10^{-5}$. These are reproducibility checks of existing controls, not independent clinical replications.

Per-patch normalization replaced a feature $x$ by $\sqrt{1536}\,x/\max(\|x\|_2,10^{-12})$ before either Hybrid or matched attention. The attention-mask candidate randomly masked six of the ten largest attention logits during training only. It was informed by attention-challenging MIL \citep{ref30} but omitted the full multi-branch architecture and therefore was not an ACMIL reproduction. Direct attention pooled unprojected 1,536-dimensional features with a linear classifier. Subspace attention instead computed attention on 256 sorted feature coordinates sampled per minibatch while pooling in the original feature space. Inference averaged logits from 21 contiguous 256-coordinate windows with stride 64, then applied sigmoid. The subspace/large-batch candidate additionally changed batch size from eight to 32 and learning rate from $5\times10^{-5}$ to $3\times10^{-4}$. This changed the update budget from 540 to 150 over 30 epochs. These subspace recipes were local adaptations informed by nnMIL \citep{ref26}, not reproductions of its complete framework; the large-batch comparison changes more than one factor.

Direct and subspace attention used 256-dimensional attention hidden layers, dropout 0.25 and no value projection. Their initial linear classifier weights were shared within each seed. The other models retained dropout 0.3. The two-bag Hybrid control averaged weighted binary cross-entropy over two independently sampled bags from each slide. The consistency variant used the same two bags and supervised loss plus $0.1\mathcal{L}_{\mathrm{cons}}$, where
\begin{equation}
\mathcal{L}_{\mathrm{cons}}=\frac{1}{2|\mathcal B|}\sum_{b\in\mathcal B}
\left[\sigma(z_{b1})-\sigma(z_{b2})\right](z_{b1}-z_{b2}).
\end{equation}
This is symmetric Bernoulli KL with gradients through both logits. It regularizes patch-bag perturbations and does not reproduce the stochastic perturbation used by dropout-only consistency methods.

Unless explicitly varied, all runs used 512 patches per bag, 30 epochs, AdamW, weight decay $10^{-4}$, batch size eight, learning rate $5\times10^{-5}$ and positive BCE weight $77/67$. The first strictly highest internal reference-bag AUROC selected the checkpoint. Final checkpoints and epoch histories were retained. There was no early stopping, scheduler, threshold search or calibration fitting. Shared first-view bags, orders and random streams followed the controlled-training convention above. Second views used the SHA-256 seed key \texttt{algorithm-v8-view2|seed|epoch|study\_slide\_id}; second-forward dropout used the original key with the suffix \texttt{dropout-view2}. Exact indices and orders were archived. The ``v8'' substring is an internal historical seed-key name and is unrelated to the manuscript version.

Eight candidates, excluding the two controls, were ranked by mean internal AUROC across the initial seeds, rounded to 12 decimal places. Ties used smaller sample standard deviation, fewer parameters and identifier order. Normalized attention was selected, and the ranking, code hashes and checkpoint hashes were saved before new external inference. Seeds 59 and 71 then evaluated normalized attention and Hybrid without changing that choice. The engineering target required internal mean AUROC gain $\geq0.01$ over Hybrid and non-negative gains in at least two initial seeds. A balanced-accuracy drop exceeding 0.02 was separately flagged. The observed AUROC gain was 0.0011 with all three differences non-negative; the AUROC target was not met. Additional seeds assess training stability on the existing data, not an independent confirmatory cohort.

\subsection{Linear controls and supplementary aggregation}

For each of the 144 training slides, raw mean features or means of normalized patch features used all available feature rows. StandardScaler and L2-regularized logistic regression were fitted on training slides only, with $C=1$, balanced class weights, LBFGS, 3,000 maximum iterations and tolerance $10^{-8}$. No hyperparameter search was performed. Evaluation used the historical 512-row reference bags; inference on all available features was also retained separately (Supplementary Table S17). These deterministic models did not enter candidate ranking.

After selection, all neural checkpoints were evaluated on historical external reference bags. Arithmetic probability ensembles retained every initial seed for each configuration and all five seeds for the selected candidate and Hybrid. For the latter two configurations, five further deterministic bags per slide used SHA-256 keys \texttt{algorithm-v8-eval|scope|study\_slide\_id|rep} for repeats 1--5. The five-bag mean was calculated separately within each seed; subsequent seed averaging gave the corresponding seed-by-bag ensemble. These analyses are distinct from the historical checkpoint's 20-bag analysis and the subspace models' feature-window logit averaging. No seed or aggregation rule was selected using external outcomes. Table S19 reports the probability-ensemble metrics; per-seed five-bag metrics and conditional 10,000-draw intervals are included in Source Data. Additional-two-seed-only probability averages are also retained there as exploratory summaries.

The full study took 620.91 seconds on one NVIDIA A800 with cached features, excluding earlier feature extraction. Timing describes this workload and is not an end-to-end deployment benchmark. Complete scores, selected epochs, run histories, configuration definitions and selection-lock records are supplied in the development-analysis archive. The study reused an external collection already examined during earlier development; its selection lock does not establish a previously untouched external test.

\subsection{Post hoc archive-tag sensitivity and descriptive Brier references}

The archive-tag sensitivity calculation excluded the 111 Source-B MALT slides tagged IHC/special stain and retained 33 H\&E-like MALT, eight HE-directory MALT and all 16 SS slides. Scores and threshold were unchanged. Table S18 reports the resulting 57-slide descriptive estimates without treating this subset as pathology-confirmed H\&E material. The 168-slide prevalence-constant Brier reference used $p=152/168$; the training-prior reference used $p=67/144$. The former uses external labels, is not deployable without that information, and was included only to contextualize the squared-error scale under marked class imbalance.
